\textbf{Step 1: an exact-margin graph family.} Write \(z=(x_2,\ldots,x_d)\) and \(D=d-1\); when \(d=1\), omit \(z\). Choose a nonincreasing \(C^\infty\) function \(\chi:[0,\infty)\to[0,1]\) equal to one on \([0,1/4]\) and zero on \([3/4,\infty)\). Put \(\psi=-2\chi\chi'\), so \(\psi\ge0\) and \(\int_0^r\psi(s)\,ds=1-\chi(r)^2\). For a smooth graph \(b\) satisfying \(\|b\|_\infty\|\psi\|_\infty<h/2\), define
\[
x_\sigma(v,z)=\sigma v+b(z)\chi(v/h)^2,
\qquad \sigma\in\{-1,1\},\quad 0\le v\le1,
\]
and define \(\tau_b\) implicitly by
\[
\tau_b(x_\sigma(v,z),z)=\sigma cv.
\]
Since
\[
\partial_vx_+=1-\frac{b(z)}h\psi(v/h)>\frac12,
\qquad
\partial_vx_-=-1-\frac{b(z)}h\psi(v/h)<-\frac12,
\]
the two branches biject \([0,1]\) onto \([b(z),1]\) and \([b(z),-1]\). Near \(v=0\), \(\tau_b=c\{x_1-b(z)\}\), while after \(v=3h/4\), \(\tau_b=cx_1\). Flatness of the cutoff makes \(\tau_b\) smooth at the joins.

Under the uniform covariate law, conditional on \(z\), the density of \(V=|\tau_b(X)|/c\) is
\[
f_{V\mid z}(v)=\frac12\left(|\partial_vx_+|+|\partial_vx_-|\right)
=\frac12\left(1-\frac b h\psi+1+\frac b h\psi\right)=1.
\]
Therefore every such graph satisfies the exact tight-margin identity
\[
P_X\{0<|\tau_b(X)|\le t\}=\frac tc,\qquad 0<t\le c.
\]
Its upper level set is \(G_b=\{(x_1,z):x_1>b(z)\}\).

Choose \(\phi\in C_c^\infty((-1/2,1/2)^D)\), with \(0\le\phi\le1\) and \(\phi=1\) on \([-1/4,1/4]^D\). Pack disjoint tangential cubes of side \(h\) with centers \(z_j\in[-1/2,1/2]^D\); their number satisfies \(M\ge c_Dh^{-D}\). Put
\[
b_\omega(z)=a h^\gamma\sum_{j=1}^M(2\omega_j-1)\phi((z-z_j)/h)^2,
\qquad \omega\in\{0,1\}^M,
\]
where \(a>0\) is fixed and sufficiently small. When \(d=1\), take \(M=1\) and \(b_\omega=a h^\gamma(2\omega_1-1)\). Because \(\gamma>d\ge1\), \(\|b_\omega\|_\infty/h\to0\). Differentiating the implicit branch maps shows that a derivative of order \(m\) of the perturbation from \(cx_1\) is bounded by \(Ca h^{\gamma-m}\), with the corresponding order-\(\gamma\) Hölder remainder bounded by \(Ca\). The linear base has norm strictly below the calibrated radius because \(c<1/4\). Thus, after fixing \(a\) small enough, every \(\tau_{b_\omega}\) belongs to the radius-one \(\gamma\)-Hölder ball. Its range is \([-c,c]\), so the near band is the whole cube and the off-band condition is vacuous. The graph lies in the interior and is crossed transversely.

\textbf{Step 2: loss separation.} For adjacent vertices, write the two local graphs as \(b_+=a\delta\phi_j^2\) and \(b_-=-a\delta\phi_j^2\), where \(\delta=h^\gamma\). On the tangential core \(Q_j\), where \(\phi_j=1\), and for \(|x_1|\le a\delta<h/4\),
\[
\tau_+(x)=c(x_1-a\delta),\qquad \tau_-(x)=c(x_1+a\delta).
\]
At every point between the two graphs, an arbitrary measurable set \(H\) disagrees with at least one truth. Hence
\[
\begin{aligned}
d_H(H,G_+;P_+)+d_H(H,G_-;P_-)
&\ge 2^{-d}c\int_{Q_j}\int_{-a\delta}^{a\delta}
\min\{x_1+a\delta,a\delta-x_1\}\,dx_1\,dz\\
&\ge c_0\delta^2h^{d-1}.
\end{aligned}
\]
This separation is uniform over every nuisance index introduced below because the target CATE and the uniform covariate density do not depend on that index.

\textbf{Step 3: the high-smoothness branch.} Let \(A\mid X\sim\operatorname{Bernoulli}(1/2)\), let \(Y\mid(A=0,X)\sim\operatorname{Bernoulli}(1/2)\), and let \(Y\mid(A=1,X)\sim\operatorname{Bernoulli}(1/2+\tau_{b_\omega}(X))\). The conditional probabilities lie in \([1/4,3/4]\); \(\pi=1/2\), \(\mu_0=1/2\), and the CATE is \(\tau_{b_\omega}\). Thus every law is in the calibrated transverse model. Adjacent CATEs differ by at most \(C\delta\) on a set of volume at most \(Ch^d\). Since \(\operatorname{KL}\{\operatorname{Bern}(p),\operatorname{Bern}(q)\}\le C(p-q)^2\) for \(p,q\in[1/4,3/4]\),
\[
\operatorname{KL}(P_\omega^{\otimes n},P_{\omega^{(j)}}^{\otimes n})\le Cn\delta^2h^d.
\]
Take \(h=(\varepsilon/n)^{1/(2\gamma+d)}\). A sufficiently small fixed \(\varepsilon\) makes adjacent total variation uniformly smaller than one. Pairing vertices across each edge and applying the preceding loss separation gives
\[
\inf_{\widehat G_n}\max_\omega\mathbb E_\omega d_H(\widehat G_n,G_\omega;P_\omega)
\ge cM\delta^2h^{d-1}\ge c\delta^2.
\]
Here the first inequality is the elementary two-point testing inequality averaged over the hypercube, and the second uses \(M\ge c_Dh^{-(d-1)}\). Since \(\delta=h^\gamma\asymp n^{-\gamma/(2\gamma+d)}=r_n^*\), this proves the high-branch lower bound.

\textbf{Step 4: the low-smoothness branch.} Put \(s=\bar s\). The strict low-branch inequality implies
\[
2s<\frac d2<\gamma.
\]
Orient an adjacent edge so that \(t_1=\tau_{-q}\ge t_0=\tau_q\), where \(q(z)=a\delta\phi_j(z)^2\), and write \(\Delta_j=t_1-t_0\). Implicit differentiation gives
\[
\partial_b\tau_b=-\frac{c\chi(v/h)^2}{1-(b/h)\psi(v/h)}
\quad\text{or}\quad
\partial_b\tau_b=-\frac{c\chi(v/h)^2}{1+(b/h)\psi(v/h)},
\]
according to the branch. Hence \(\Delta_j\ge0\), and the divided-difference identity gives
\[
\Delta_j(x,z)=\delta\phi_j(z)^2A_j(x,z),
\qquad
A_j=-2a\int_0^1\left.\partial_b\tau_b\right|_{b=(2r-1)q}\,dr.
\]
The function \(A_j\) is nonnegative, smooth, supported in a region of volume \(O(h^d)\), and satisfies \(|D^mA_j|\le C_mh^{-m}\).

Set \(w=\delta^{1/(2s)}=h^{\gamma/(2s)}\), so \(w<h\). A concrete smooth squared partition of unity on the support of \(A_j\phi_j^2\) is obtained as follows. Start with a nonnegative bump \(\zeta_0\) supported in \((-1,1)^d\) whose integer translates cover \(\mathbb R^d\), and put
\[
\zeta(u)=\frac{\zeta_0(u)}{\{\sum_{r\in\mathbb Z^d}\zeta_0(u-r)^2\}^{1/2}}.
\]
After translating and rescaling by \(w\), the functions \(\chi_l\) satisfy \(\sum_l\chi_l^2=1\), have overlap bounded only by \(d\), have support volume \(O(w^d)\), and obey \(|D^m\chi_l|\le C_mw^{-m}\). Only \(K\le Ch^dw^{-d}\) of them meet the coarse support.

Choose a fixed small scale \(a_*\) and independent Rademacher signs \(\lambda_l\). Define
\[
u_\lambda=a_*w^\alpha\sum_l\lambda_l\phi_j\chi_l,
\qquad
v_\lambda=\frac{w^\beta}{4a_*}\sum_l\lambda_lA_j\phi_j\chi_l.
\]
The graph amplitude \(a\) is chosen small relative to \(a_*^2\), so both nuisance Hölder norms fit in the calibrated radius-one balls. Since \(w^{\alpha+\beta}=w^{2s}=\delta\),
\[
\mathbb E_\lambda u_\lambda=\mathbb E_\lambda v_\lambda=0,
\qquad
4\mathbb E_\lambda(u_\lambda v_\lambda)=\Delta_j.
\]
The product rule, \(w<h\), and the derivative bounds on \(A_j,\phi_j,\chi_l\) show directly that \(u_\lambda\) has uniformly bounded \(\alpha\)-Hölder norm and \(v_\lambda\) has uniformly bounded \(\beta\)-Hölder norm.

If \(\alpha\ge\beta\), define the two conditional four-cell laws by
\[
\begin{aligned}
p_\lambda(a_0,y_0\mid x)
&=\frac14+(y_0-\tfrac12)v_\lambda+(2a_0-1)(2y_0-1)\frac{t_1}{4},\\
q_\lambda(a_0,y_0\mid x)
&=\frac14+(a_0-\tfrac12)u_\lambda+(y_0-\tfrac12)(v_\lambda+t_0u_\lambda)
 +(2a_0-1)(2y_0-1)\left(\frac{t_0}{4}+u_\lambda v_\lambda\right).
\end{aligned}
\]
Summing over \(y_0\) and dividing the \(y_0=1\) cell by the treatment marginal shows that \(p_\lambda\) has \((\pi,\mu_0,\tau)=(1/2,1/2+v_\lambda-t_1/2,t_1)\), whereas \(q_\lambda\) has \((\pi,\mu_0,\tau)=(1/2+u_\lambda,1/2+v_\lambda-t_0/2,t_0)\). If \(\beta\ge\alpha\), use the same \(q_\lambda\) and instead put
\[
p_\lambda(a_0,y_0\mid x)=\frac14+(a_0-\tfrac12)u_\lambda+(y_0-\tfrac12)t_1u_\lambda+(2a_0-1)(2y_0-1)\frac{t_1}{4}.
\]
Then \(p_\lambda\) has \((\pi,\mu_0,\tau)=(1/2+u_\lambda,1/2-t_1/2,t_1)\), while \(q_\lambda\) has the same parameters as before. In either case the one-observation mixtures agree exactly because \(4\mathbb E(u_\lambda v_\lambda)=t_1-t_0\). Since \(\max(\alpha,\beta)\le\alpha+\beta=2s<\gamma\), the terms involving \(t_i\) satisfy the required nuisance smoothness. Small fixed \(a,a_*\) keep all four cell probabilities bounded away from zero, give strict overlap, and place \(\pi\) and \(\mu_0\) in their radius-one Hölder balls. The exact-margin graph calculation is unaffected. Pasting independent copies over the disjoint outer graph cells therefore produces laws in \(\mathcal M_{\mathrm{tr}}(\gamma;\kappa)\).

The bounded-overlap mixture lemma applies with amplitudes of orders \(w^\alpha\) and \(w^\beta\). Because \(\max(\alpha,\beta)\ge s\),
\[
\min\{w^{4\alpha},w^{4\beta}\}\le w^{4s}=\delta^2,
\]
and consequently
\[
H^2(\overline P_\omega^{\otimes n},\overline P_{\omega^{(j)}}^{\otimes n})
\le Cn^2Kw^{2d}\delta^2
\le Cn^2h^dw^d\delta^2
=Cn^2h^{2\gamma+d+d\gamma/(2s)}.
\]
Let \(T=1+d/(4s)+d/(2\gamma)\) and choose \(h=(\varepsilon/n)^{1/(\gamma T)}\). The identity
\[
2\gamma+d+\frac{d\gamma}{2s}=2\gamma T
\]
makes the last bound at most \(C\varepsilon^2\). Moreover, the strict low-branch inequality is equivalent to \(d/(2s)>T\), and hence
\[
nw^d=nh^{d\gamma/(2s)}\longrightarrow0,
\]
which verifies the small-collision hypothesis of the mixture lemma. Choose \(\varepsilon\) so that adjacent mixture total variation is bounded away from one. Pairing adjacent fuzzy mixtures and using the uniform loss separation gives
\[
\inf_{\widehat G_n}\sup_{P\in\mathcal M_{\mathrm{tr}}(\gamma;\kappa)}
\mathbb E_P d_H(\widehat G_n,G_P;P)
\ge cM\delta^2h^{d-1}\ge c\delta^2.
\]
Finally, \(\delta=h^\gamma\asymp n^{-1/T}=r_n^*\). For \(d=1\), \(M=1\), and the same argument is the corresponding direct or fuzzy two-hypothesis experiment. This completes both branches.